% ══════════════════════════════════════════════════════════════════════════════
%  EXPERIMENT 4 — MECHANICAL ADVANTAGE               (Lesson 10, T3 Week 4)
%  From Lesson 10 of the A4 Forces Lesson Booklet; answer keys 10.* kept.
% ══════════════════════════════════════════════════════════════════════════════
\experimentdiv{Mechanical Advantage}{Balancing a lever \textbullet{} Effort arm and load arm \textbullet{} Two ways to find the MA}{10}\label{sec:4}
\lessoncodes{VC2S8U13 \textbullet{} VC2S8I01, I02, I03, I04, I05, I06}

\vocabhead
\begin{ctab}
\begin{tabular}{|SW{4.2cm}|W{11.6cm}|}
  \hline\hd & \thead{Word} & \thead{Definition (copy from the slide)} \\ \hline
  & \scaf{\textbf{Steelyard}} & \scaf{An old weighing machine: a lever with a sliding weight that balances the load.} \\ \hline
  & \textbf{Pivot} & \cell{10.v2} \\ \hline
  & \textbf{Effort arm} & \cell{10.v3} \\ \hline
  & \textbf{Load arm} & \cell{10.v4} \\ \hline
  & \textbf{Balanced lever} & \cell{10.v5} \\ \hline
\end{tabular}
\end{ctab}

\begin{lipanel}
\inquiryli{Measure distances carefully and check one calculation against another.}
\understandli{Use the law of the lever to find a lever's mechanical advantage.}
\sctitle
\begin{itemize}[leftmargin=1.4em, itemsep=0pt, topsep=1pt]
  \item I can measure each arm from the fulcrum to the centre of the mass.
  \item I can find the MA from the arms and from the forces.
  \item I can draw a smooth curve of best fit when the trend is not a straight line.
\end{itemize}
\end{lipanel}

\begin{skillsfocus}
\focus{QP1}{Write a hypothesis in the faculty frame.}
\focus{PC2}{Name the variables, with two controlled.}
\focus{PC6}{Measure each arm to the centre of the mass, three times.}
\focus{PA1}{Record trials, means and MAs in a THUD table.}
\focus{PA2}{Plot the means and draw a curve of best fit.}
\focus{PA3}{Calculate the MA two ways.}
\focus{EV2}{Compare your results with the law of the lever.}
\end{skillsfocus}

% ─── READING ─────────────────────────────────────────────────────────────────
\exppart{Before the lab}{Teal}

\begin{readingbox}{Reading --- The Steelyard: Weighing with a Lever}
Long before electronic scales, market sellers weighed goods with a \emph{steelyard}. A steelyard is a lever: a long bar hanging from a hook near one end, which acts as the pivot. The goods hang from the short side, close to the pivot. A small, heavy weight slides along the long side, which is marked with a scale. The seller slides the weight until the bar hangs level, then reads the weight of the goods off the scale.

How can a small weight balance a heavy load? The answer is distance. A lever balances when effort $\times$ effort arm equals load $\times$ load arm. Far from the pivot, a small weight turns the bar as strongly as a large weight close to it. A sliding weight ten times further from the pivot than the goods can balance goods ten times heavier than itself.

Steelyards are very old: bronze ones have been found at Pompeii, the Roman town buried by Vesuvius in AD~79. About three hundred years earlier, the Greek mathematician Archimedes had worked out the law of the lever that makes them work.

In this experiment you will do what a steelyard does: balance a heavy load near the pivot with a lighter mass further away, measure both arms, and work out your lever's mechanical advantage.
\end{readingbox}

\rcreset
\rcqk{10.c1}{Describe how a market seller uses a steelyard to find the weight of some goods.}{2.4cm}
\rcqk{10.c2}{Explain how a small sliding weight can balance a much heavier load. Use the words \emph{pivot} and \emph{distance}.}{2.4cm}
\rcqk{10.c3}{The sliding weight sits ten times further from the pivot than the goods. How heavy can the goods be, compared with the sliding weight?}{2.2cm}
\rcqk{10.c4}{In this experiment, what will you balance, and what will you measure?}{2.2cm}

% ─── TECHNIQUE ────────────────────────────────────────────────────────────────
\exppart{Technique spotlight}{Blue}

\begin{technique}{Measuring an arm}{PC6}
Measure each arm from the fulcrum to the \emph{centre} of the mass, not its edge. If the stick balances a little to one side of 50\,cm on its own, use that point as the fulcrum and measure every arm from there.
\end{technique}

\begin{infobox}{Mechanical advantage, two ways}
\textbf{From the arms:} \quad $\text{MA} = \text{effort arm} \div \text{load arm}$\\[2pt]
\textbf{From the forces:} \quad $\text{MA} = \text{load} \div \text{effort}$\\[4pt]
Weight is proportional to mass, so you can compare the forces using the masses in grams. When the lever is balanced, the two answers should agree.
\end{infobox}

\begin{workedex}{Mechanical advantage of a balanced lever}
A 100\,g load stands 20\,cm from the fulcrum. A 50\,g effort balances it 40\,cm from the fulcrum.

\smallskip
\textbf{From the arms:}\quad $40 \div 20 = \mathbf{2}$ \qquad \textbf{From the forces:}\quad $100 \div 50 = \mathbf{2}$

\smallskip
They agree. This is the law of the lever: $100 \times 20 = 50 \times 40 = 2000$ on both sides.
\end{workedex}

\begin{wedo}{Mechanical advantage of a balanced lever}
A 150\,g load stands 10\,cm from the fulcrum. A 50\,g effort balances it 30\,cm from the fulcrum.
\pt{a}Find the MA from the arms.
\wbox{10.wd1a}{1.8cm}
\pt{b}Find the MA from the forces.
\wbox{10.wd1b}{1.8cm}
\pt{c}Show that the law of the lever works for this lever.
\wbox{10.wd1c}{1.8cm}
\end{wedo}

\begin{technique}{When the trend is a curve}{PA2}
Not every trend is a straight line. If your points bend, draw one smooth \textbf{curve of best fit} through the middle of them --- still never dot to dot, and never a ruler line through a curve.
\end{technique}

% ─── PLAN, SAFETY AND METHOD ─────────────────────────────────────────────────
\exppart{Plan, safety and method}{Teal}

\textbf{Aim:} To investigate the effect of changing the effort mass on the effort arm needed to balance a lever.

\task Complete the variables table. \skr{PC2}
\begin{ctab}
\begin{tabular}{|SW{4.4cm}|W{10.2cm}|}
  \hline\hd & \thead{Variable} & \thead{In this experiment} \\ \hline
  & \scaf{Independent} & \scaf{The effort mass: 50, 100, 150 and 200\,g} \\ \hline
  & Dependent & \cell{10.t1-r2} \\ \hline
  & Controlled & \cell{10.t1-r3} \\ \hline
  & Controlled & \cell{10.t1-r4} \\ \hline
\end{tabular}
\end{ctab}

\needspace{8\baselineskip}
\task Write your hypothesis. \skr{QP1}

\scaf{If the effort mass is increased, the effort arm needed to balance the lever will \dots\ because \dots}
\abox{10.pred}{2.5cm}

\subsubsection*{Equipment}
Metre stick \textbullet{} triangular prism or thick pencil (fulcrum) \textbullet{} one 100\,g mass (load) \textbullet{} four 50\,g slotted masses (effort) \textbullet{} sticky tack

\begin{center}
\diagExpFourSetup
\expfigcap{Measure each arm from the fulcrum to the centre of the mass. Here a 50\,g effort balances the 100\,g load at an effort arm of 40\,cm.}
\end{center}

\needspace{12\baselineskip}
\task Complete the risk assessment. \skr{PC3}
\begin{ctab}
\begin{tabular}{|SW{4.0cm}|W{5.2cm}|W{5.4cm}|}
  \hline\hd & \thead{Risk (what could go wrong)} & \thead{Prevention} & \thead{Responding (if it happens)} \\ \hline
  & \scaf{The stick tips and drops the masses on a foot} & \scaf{Hold the stick while adding or moving a mass; keep feet clear of the bench edge} & \scaf{Tell the teacher; cold pack for a bruise} \\ \hline
  & Masses roll off the bench & & \\ \hline
\end{tabular}
\end{ctab}

\subsubsection*{Method}
\begin{enumerate}[itemsep=2pt]
  \item Rest the metre stick on the fulcrum at the 50\,cm mark and check that it balances level.
  \item Stand the 100\,g load on the 30\,cm mark, so the load arm is 20\,cm. Fix it with a little sticky tack. It stays there for the whole experiment.
  \item Stand one 50\,g mass on the other side of the fulcrum. Slide it until the stick balances level. Record the effort arm in Table 4.1.
  \item Take the effort mass off and repeat step 3 twice more, so you have three trials.
  \item Repeat steps 3--4 with an effort of 100\,g, 150\,g and 200\,g (stack the 50\,g masses).
\end{enumerate}

% ─── RESULTS ─────────────────────────────────────────────────────────────────
\exppart{Results}{Amber}

\subsubsection*{Table 4.1: Effort arm and MA for each effort mass (load 100\,g, load arm 20\,cm) \skr{PA1}}
\begin{ctab}
\begin{tabular}{|SZ{1.9cm}|Z{1.5cm}|Z{1.5cm}|Z{1.5cm}|Z{2.2cm}|Z{2.5cm}|Z{2.5cm}|}
  \hline\hd & \thead{Effort mass (g)} & \thead{Effort arm, trial 1 (cm)} & \thead{Effort arm, trial 2 (cm)} & \thead{Effort arm, trial 3 (cm)} & \thead{Mean effort arm (cm)} & \thead{MA = effort arm $\div$ load arm} & \thead{MA = load $\div$ effort} \\ \hline
  & \scaf{e.g.\ 50} & \scaf{39} & \scaf{41} & \scaf{40} & \scaf{40} & \scaf{$40 \div 20 = 2$} & \scaf{$100 \div 50 = 2$} \\ \hline
  & 50  & \cell{10.res-r2c2} & \cell{10.res-r2c3} & \cell{10.res-r2c4} & \cell{10.res-r2c5} & \cell{10.res-r2c6} & \cell{10.res-r2c7} \\ \hline
  & 100 & \cell{10.res-r3c2} & \cell{10.res-r3c3} & \cell{10.res-r3c4} & \cell{10.res-r3c5} & \cell{10.res-r3c6} & \cell{10.res-r3c7} \\ \hline
  & 150 & \cell{10.res-r4c2} & \cell{10.res-r4c3} & \cell{10.res-r4c4} & \cell{10.res-r4c5} & \cell{10.res-r4c6} & \cell{10.res-r4c7} \\ \hline
  & 200 & \cell{10.res-r5c2} & \cell{10.res-r5c3} & \cell{10.res-r5c4} & \cell{10.res-r5c5} & \cell{10.res-r5c6} & \cell{10.res-r5c7} \\ \hline
\end{tabular}
\end{ctab}

\needspace{20\baselineskip}
\task Draw a line graph of your results. Put \textbf{effort mass (g)} on the horizontal axis and \textbf{mean effort arm (cm)} on the vertical axis. Plot each point, then draw a smooth curve of best fit. Give your graph a title. \skr{PA2}
\graphgrid{9cm}

% ─── ANALYSIS ────────────────────────────────────────────────────────────────
\exppart{Analysis and evaluation}{Coral}

\task What happened to the effort arm as the effort mass increased? What happened to the MA? Use numbers from your table. \skr{PA5}
\abox{10.t2}{3.5cm}

\needspace{10\baselineskip}
\task Compare your two MA columns. Do they agree? Suggest one reason why they might not agree exactly. \skr{EV2}
\abox{10.t3}{3.5cm}

\needspace{8\baselineskip}
\task Another group balanced the 100\,g load (load arm 20\,cm) with an 80\,g effort. Use the law of the lever to predict their effort arm. \skr{PA3}
\wbox{10.t4}{2.6cm}

\needspace{8\baselineskip}
\task A group's effort arms for 100\,g were 19, 21 and 27\,cm. Which value looks wrong? Find the mean of the other two.
\wbox{10.t5}{2.6cm}

\needspace{10\baselineskip}
\task Use the law of the lever to explain \emph{why} a heavier effort mass balanced the load closer to the fulcrum.
\abox{10.t6}{4cm}

\needspace{12\baselineskip}
\task Write a conclusion: state the trend, back it with numbers from your table, and say whether your hypothesis was supported.

\scaf{As the effort mass increased from 50\,g to 200\,g, the mean effort arm decreased from \dots\,cm to \dots\,cm. This supports / does not support the hypothesis.}
\abox{10.concl}{5cm}

\begin{yr8box}
\begin{keepext}
\exttask For each row of your results, calculate effort mass $\times$ mean effort arm. What do you notice? Use the law of the lever to explain why this number should be the same every time, and what it should equal.
\wbox{10.x8}{4.5cm}
\end{keepext}
\begin{keepext}
\exttask A trebuchet is a lever too, but it is built with an MA \emph{less} than 1. Describe how its effort and load are placed differently from this experiment, and explain why it is built that way.
\abox{10.x9}{4.5cm}
\end{keepext}
\end{yr8box}

% ─── SKILL CHECK ─────────────────────────────────────────────────────────────
\exppart{Skill check}{NavyBlue}

\begin{skillcheck}
\scrow{QP1}{My hypothesis names the effort mass and the effort arm.}
\scrow{PC2}{I named the variables, with two controlled.}
\scrow{PC6}{I measured each arm to the centre of the mass.}
\scrow{PA1}{My table has a title, headings, units in the headings and matching decimals.}
\scrow{PA2}{I drew a smooth curve of best fit.}
\scrow{PA3}{I found the MA two ways.}
\scrow{EV2}{I compared my results with the law of the lever.}
\end{skillcheck}

\task Which skill will you focus on in Experiment 5? Explain what you will do differently.
\abox{10.sc}{5cm}
